\chapter{학부모 관리}

\begin{summarybox}{한눈에}
본 기관 학생의 학부모 계정을 관리하고, 학생-학부모 \textbf{연결(매칭)} 을 승인합니다. 학부모는 자녀의 학습 현황과 통합 성적표를 자기 화면에서 볼 수 있어요.
\end{summarybox}

\kfShot{05-parents}{학부모 관리 — 본 기관 학생-학부모 연결 목록}

\section*{화면 구조}

\begin{screenbox}{학부모 ↔ 학생 연결 목록 + 신청 처리}
\noindent
\begin{tabular}{@{}p{2.5cm}p{2.5cm}p{2.5cm}p{2cm}p{2cm}@{}}
\toprule
\textbf{학부모} & \textbf{자녀(학생)} & \textbf{관계} & \textbf{상태} & \textbf{액션} \\
\midrule
김부모 (010-...) & 홍길동 & 부 & active & \inacttab{해제} \\
이부모 (010-...) & 김영희 & 모 & pending & \activetab{승인} \inacttab{거절} \\
\bottomrule
\end{tabular}
\end{screenbox}

\section*{여기서 할 수 있는 일}
\begin{itemize}[leftmargin=1.2em]
  \item 학부모-학생 \textbf{연결 신청 승인/거절}
  \item 기존 연결 \textbf{직접 추가/해제}
  \item 학부모 연락처 확인 (학생 학부모 연락처는 학생 상세에서도 확인 가능)
\end{itemize}

\section*{자주 쓰는 흐름}
\begin{enumerate}
  \item 학부모가 회원가입 시 ``자녀 학생 아이디''를 입력 → 신청 발송됨
  \item 기관 관리자가 ``학부모 연결 대기'' 카드(대시보드) → 이 화면 → 승인
  \item 승인 후 학부모는 자녀 통합 성적표·학습 현황 조회 가능
\end{enumerate}

\begin{tipbox}{알아두면 좋은 점}
\begin{itemize}[leftmargin=1.2em]
  \item 본 기관 학생의 학부모 연결만 보입니다 — 다른 기관 학생은 자동 필터됨
  \item 학부모 한 명이 여러 자녀 연결 가능 (각각 승인)
  \item 연결 해제는 학생 본인이 학부모 화면에서도 할 수 있어요 (양방향 권한)
\end{itemize}
\end{tipbox}
